\documentclass[10pt,a4paper]{amsart}
\usepackage[margin=19mm]{geometry}
\usepackage{amsmath,amssymb,mathtools}
\usepackage[scr=rsfs,cal=euler]{mathalfa}
\usepackage{xfrac}
\usepackage[unicode,hidelinks,bookmarks=false]{hyperref}
\newcommand{\ie}{i.e.\ }
\newcommand{\cf}{cf.\ }
\newcommand{\resp}{respectively\ }
\newcommand{\Subsection}{\subsection}
\providecommand{\nobreakdash}{\nobreak\hbox{-}}
\setlength{\emergencystretch}{2em}
\multlinegap0pt
\usepackage{xcolor}
\usepackage{amssymb}
\usepackage{mathtools}
\usepackage{dsfont}
\usepackage{subfig}
\usepackage{float}
\usepackage{enumerate}
\usepackage[shortlabels]{enumitem}
\usepackage{bm}
\usepackage{tikz}
\usepackage{multirow}
\usepackage{caption}
\newtheorem{lemma}{Lemma}
\newtheorem{proposition}{Proposition}
\newcommand{\SL}{\mathrm{SL}}
\newcommand{\bZ}{\mathbb{Z}}
\newcommand{\cM}{\mathcal{M}}
\title{Inequalities involving polynomials and quasimodular forms}
\author{Seewoo Lee}
\date{Source: arXiv:2602.10536v1 (CC BY 4.0)}
\begin{document}
\maketitle

\noindent\emph{Source and licence.} Inequalities involving polynomials and quasimodular forms. Version arXiv:2602.10536v1 (CC BY 4.0), distributed by arXiv under the Creative Commons Attribution 4.0 International licence, http://creativecommons.org/licenses/by/4.0/, as recorded in the arXiv licence metadata for this exact version and shown on the abstract page https://arxiv.org/abs/2602.10536v1. Full licence notice: \texttt{licenses/arxiv-2602.10536-CC-BY-4.0.txt}.\par\smallskip

\noindent\textit{Editorial scope.} Counted unit: the article's proposition prop:Xw1limit (source0.tex lines 1099--1108). The article's complete original proof of it is delivered with it (source0.tex lines 1109--1149, one \texttt{proof} environment of 41 lines). No mathematics is written, repaired or re-derived: the statement and the proof are the article's own lines, in the article's own notation, and no retained line is substituted or re-flowed.  Retained uncounted as the source-local context the proof needs: lines 729--744.  Nothing was omitted from any retained span, so the package carries no omission record.  No retained line contains a citation, so the wrapper carries no bibliography and no bibliography entry.  No retained line contains a figure environment, an image-inclusion command or drawing code, so no image is delivered and the package declares no asset dependency.  Editorial formatting only, in addition: an AMS \texttt{amsart} preamble that restates the article's own declarations and macros used by the retained lines, namely the environments \texttt{lemma}, \texttt{proposition} on separate counters, together with the standard packages those lines need.  The retained environments print in this document as Lemma 1, Proposition 1 in order of appearance, whereas the article prints the same items with the numbering of its own full document.  The complete native source of the article as distributed in the arXiv e-print for this exact version is bundled unchanged as \texttt{sources/194408/source0.tex}.
\begin{lemma}
    \label{lem:n0coeff}
    For $w \equiv 0 \pmod{6}$, the constant terms $\alpha_{w,0}$ and $\beta_{w-2,0}$ satisfy
    \begin{align}
        \alpha_{w+2,0} &= -\frac{w-1}{w+1} \alpha_{w,0} \label{eqn:alphaw2} \\
        \alpha_{w+4,0} &= \alpha_{w,0} \label{eqn:alphaw4} \\
        \alpha_{w+6,0} &= -\frac{w(w+6)}{432(w+1)(w+5)} \alpha_{w,0} \label{eqn:alphaw6}
    \end{align}
    and $\beta_{w,0} = -\alpha_{w+2,0}$ for all $w \ge 4$.
    In particular, $A_{w}$ and $B_{w-2}$ are not cusp forms for all $w \ge 6$.
    Also, for $w \equiv 0 \pmod{6}$, we have
    \begin{equation}
        \label{eqn:alphawclosed}
        \alpha_{w,0} = (-1)^{\frac{w}{6}}\frac{(w/6)! (w/3)! (w/2)! }{2w \cdot w!}.
    \end{equation}
\end{lemma}

\begin{proposition}
    \label{prop:Xw1limit}
    Let $X_{w, 1}$ be the normalized extremal quasimodular form of weight $w$ and depth $1$.
    Then
    \begin{align}
        \lim_{t \to 0^+} \frac{X_{w, 1}(it)}{t X_{w, 1}'(it)} &= \frac{2\pi}{w-1} \label{eqn:extremallimit} \\
        \lim_{t \to 0^+} t^{w-1} X_{w, 1}(it) &= -\frac{6(-1)^{\frac{w}{2}}\beta_{w-2}}{\pi} \label{eqn:extremallimit2}
    \end{align}
    for all $w \ge 6$.
\end{proposition}

\begin{proof}
    The transformation law gives
    \[
        X_{w, 1}\left(-\frac{1}{z}\right) = z^w A_w(z) + \left(z^2 E_2(z) - \frac{6iz}{\pi}\right) z^{w-2}B_{w-2}(z) = z^w X_{w, 1}(z) - \frac{6iz^{w-1}}{\pi} B_{w-2}(z)
    \]
    and $z = it$ gives
    \begin{equation}
        X_{w,1} \left(\frac{i}{t}\right) = (-1)^{\frac{w}{2}}\left[t^w X_{w, 1}(it) - \frac{6 t^{w-1}}{\pi} B_{w-2}(it)\right]. \label{eqn:Xw1trans}
    \end{equation}
    By differentiating $X_{w, 1}$ and rewriting in terms of Serre derivatives, we have
    \begin{align*}
        X_{w, 1}' &= A_{w}' + E_2' B_{w-2} + E_2 B_{w-2}' \\
        &= \partial_w A_w + \frac{w}{12} E_2 A_{w} + \frac{E_2^2 - E_4}{12} B_{w-2} + E_2 \left(\partial_{w-2} B_{w-2} + \frac{w-2}{12} E_2 B_{w-2}\right) \\
        &= \left(\partial_{w} A_{w} - \frac{1}{12} E_4 B_{w-2}\right) + E_2 \left(\frac{w}{12} A_w + \partial_{w-2} B_{w-2}\right) + E_2^2 \cdot \frac{w-1}{12} B_{w-2} \\
        &=: \widetilde{A}_{w+2} + E_2 \widetilde{B}_{w} + E_2^2 \widetilde{C}_{w-2},
    \end{align*}
    where
    \begin{align*}
        \widetilde{A}_{w+2} &:= \partial_{w} A_{w} - \frac{1}{12} E_4 B_{w-2} \in \cM_{w+2}(\SL_2(\bZ)), \\
        \widetilde{B}_{w} &:= \frac{w}{12} A_w + \partial_{w-2} B_{w-2} \in \cM_{w}(\SL_2(\bZ)), \\
        \widetilde{C}_{w-2} &:= \frac{w-1}{12} B_{w-2} \in \cM_{w-2}(\SL_2(\bZ)).
    \end{align*}
    Then applying the transformation law with $z = it$ gives
    \begin{align}
        X_{w, 1}'\left(\frac{i}{t}\right) &= (-1)^{\frac{w}{2} + 1} \left[t^{w+2} X_{w, 1}'(it) + \frac{t^{w+1}}{\pi} (-6 \widetilde{B}_{w}(it) - (w-1) B_{w-2}(it) E_{2}(it)) + \frac{t^w}{\pi^2} \cdot 3(w-1) B_{w-2}(it)\right] \nonumber\\
        &= (-1)^{\frac{w}{2} + 1} \left[t^{w+2} X_{w, 1}'(it) + \frac{t^{w+1}}{\pi} \left(-\frac{w}{2} X_{w, 1}(it) - 6 B_{w}'(it)\right) + \frac{t^w}{\pi^2} \cdot 3(w-1) B_{w-2}(it)\right]  \label{eqn:dXw1trans}
    \end{align}
    By combining \eqref{eqn:Xw1trans} and \eqref{eqn:dXw1trans}, we have
    \begin{align*}
        \lim_{t \to 0^+} \frac{X_{w, 1}(it)}{t X_{w, 1}'(it)} &= \lim_{t \to \infty} \frac{X_{w, 1}(i/t)}{\frac{1}{t} X_{w, 1}'(i/t)} \\
        &= \lim_{t \to \infty} -\frac{t^w X_{w, 1}(it) - \frac{6 t^{w-1}}{\pi} B_{w-2}(it)}{t^{w+1} X_{w, 1}'(it) + \frac{t^w}{\pi} \left(-\frac{w}{2} X_{w, 1}(it) - 6 B_{w-2}'(it)\right) + \frac{t^{w-1}}{\pi^2} \cdot 3(w-1) B_{w-2}(it)} \\
        &= -\lim_{t \to \infty} \frac{t X_{w, 1}(it) - \frac{6}{\pi} B_{w-2}(it)}{t^2 X_{w, 1}'(it) + \frac{t}{\pi} \left(-\frac{w}{2} X_{w, 1}(it) - 6 B_{w-2}'(it)\right) + \frac{1}{\pi^2} \cdot 3(w-1) B_{w-2}(it)}.
    \end{align*}
    Since $X_{w, 1}, X_{w, 1}'$ and $-\frac{w}{2} X_{w, 1} - 6 B_{w-2}'$ are cusp forms, the limits of the terms involving them vanish as $t \to \infty$.
    By Lemma \ref{lem:n0coeff}, $B_{w-2}$ is not a cusp form and the final limit equals to
    \[
    -\frac{-\frac{6}{\pi} \cdot \beta_{w-2}}{\frac{1}{\pi^2} \cdot 3(w-1) \cdot \beta_{w-2}} = \frac{2\pi}{w-1}
    \]
    which proves \eqref{eqn:extremallimit}.
    The limit \eqref{eqn:extremallimit2} follows from \eqref{eqn:Xw1trans} and the cuspidality of $X_{w, 1}$.
\end{proof}

\end{document}
